\section{Exact semantic error and the no-log accelerated upper bound}
\label{sec:signed-star-semantic-upper}

We now prove the central result.  The proof has two ingredients: an exact
critically damped residual recurrence and an exact conversion of its two
signed layers into semantic error.

\begin{lemma}[Exact block recurrence]
\label{lem:signed-star-block-recurrence}
Let $x_j$ denote the aggregate residual processed by block $j$, with the
initial center block indexed by $j=0$.  Then
\begin{equation}
  x_0=\gamma_\alpha,
  \qquad
  x_1=2\lambda\gamma_\alpha,
  \qquad
  x_{j+1}=2\lambda x_j-\lambda^2x_{j-1},
  \label{eq:signed-star-block-recurrence}
\end{equation}
and hence
\begin{equation}
  x_j=\gamma_\alpha(j+1)\lambda^j.
  \label{eq:signed-star-block-solution}
\end{equation}
After $k\ge1$ completed blocks, the center and aggregate-leaf residuals, in
an order determined by the parity of $k$, are
\begin{equation}
  a_k:=x_k=\gamma_\alpha(k+1)\lambda^k,
  \qquad
  -b_k:=-\lambda^2x_{k-1}
      =-\gamma_\alpha k\lambda^{k+1}.
  \label{eq:signed-star-two-layers}
\end{equation}
In particular $a_k>b_k\ge0$.
\end{lemma}

\begin{proof}
Immediately before block $j$, the layer being processed contains $x_j$ and
the previously processed layer contains $-\lambda^2x_{j-1}$.  The block
reflects the former to $-\lambda^2x_j$ and adds $2\lambda x_j$ to the latter.
Thus the next processed residual is
$x_{j+1}=2\lambda x_j-\lambda^2x_{j-1}$.  Its characteristic polynomial is
$(z-\lambda)^2$.  The initial values determine
$x_j=\gamma_\alpha(j+1)\lambda^j$.  The post-block description follows
directly, and $a_k/b_k=(k+1)/(k\lambda)>1$.
\end{proof}

\begin{lemma}[Exact semantic-error identity]
\label{lem:signed-star-semantic-identity}
Let $p^{(k)}$ be the iterate after $k\ge0$ complete blocks.  Then
\begin{equation}
  \boxed{
  \mathcal E(p^{(k)})
  =\frac{\lambda^k}{2m}
      \left(1+t^2+2tk\right),
  \qquad t=\sqrt\alpha .}
  \label{eq:signed-star-exact-semantic-error}
\end{equation}
\end{lemma}

\begin{proof}
Let $e=\pi-p^{(k)}=H^{-1}r$, and let
$E_L:=\sum_{\ell\in L}e_\ell$.  For any symmetric star residual pair
$(R,T)$, the two equations $He=r$ reduce to
\begin{equation}
  e_o-c_\alpha E_L=R,
  \qquad
  E_L-c_\alpha e_o=T.
  \label{eq:signed-star-error-system}
\end{equation}
Therefore
\begin{equation}
  e_o=\frac{R+c_\alpha T}{1-c_\alpha^2},
  \qquad
  E_L=\frac{T+c_\alpha R}{1-c_\alpha^2}.
  \label{eq:signed-star-error-solution}
\end{equation}
Every leaf has error $E_L/m$, while the center has degree $m$, so
\begin{equation}
  \mathcal E(p^{(k)})
  =\frac{1}{m(1-c_\alpha^2)}
    \max\{|R+c_\alpha T|,|T+c_\alpha R|\}.
  \label{eq:signed-star-semantic-max}
\end{equation}

For $k\ge1$, the unordered residual pair is $\{a_k,-b_k\}$ from
\eqref{eq:signed-star-two-layers}.  Whenever $a\ge b\ge0$ and
$0\le c\le1$,
\[
  \max\{|a-cb|,|ca-b|\}=a-cb.
\]
Indeed, compare $a-cb$ with $ca-b$ if the latter is nonnegative, and with
$b-ca$ otherwise.  Hence
\begin{align}
  \mathcal E(p^{(k)})
  &=\frac{a_k-c_\alpha b_k}{m(1-c_\alpha^2)} \notag\\
  &=\frac{\lambda^k\bigl(1+k(1-c_\alpha\lambda)\bigr)}
          {m(1+c_\alpha)}.
  \label{eq:signed-star-semantic-before-cancel}
\end{align}
Now
\begin{equation}
  c_\alpha=\frac{1-t^2}{1+t^2},
  \qquad
  1+c_\alpha=\frac{2}{1+t^2},
  \qquad
  1-c_\alpha\lambda=\frac{2t}{1+t^2}.
  \label{eq:signed-star-cancellation-identities}
\end{equation}
Substitution proves \eqref{eq:signed-star-exact-semantic-error}.  At $k=0$
the formula gives $(1+t^2)/(2m)=\pi_o/m$, which is the initial semantic
error, so the identity also holds there.
\end{proof}

The last identity in \eqref{eq:signed-star-cancellation-identities} is the
essential cancellation.  The residual magnitude contains $k\lambda^k$, but
the semantic combination contains only $tk\lambda^k$.  Thus the critical
double-root transient is attenuated by exactly the spectral scale
$t=\sqrt\alpha$.

\begin{theorem}[Predetermined-stop accelerated SOR on the star]
\label{thm:signed-star-semantic-upper}
Let $0<\alpha<1$, let
$0<\varepsilon_{\rm ppr}\le1/16$, and choose $m$ as in
\eqref{eq:signed-star-size}.  The signed-SOR algorithm of
\Cref{sec:signed-star-algorithm}, stopped after
$k_\star=\lceil2/\sqrt\alpha\rceil$ blocks, satisfies
\begin{equation}
  \mathcal E(p^{(k_\star)})\le\varepsilon_{\rm ppr}.
  \label{eq:signed-star-upper-accuracy}
\end{equation}
Including the sparse seed read and $m+1$ explicit output writes, its work is
\begin{equation}
  W_{\rm SOR}
  \le m\left\lceil\frac{2}{\sqrt\alpha}\right\rceil+m+2
  =O\!\left(
      \frac{1}{\sqrt\alpha\,\varepsilon_{\rm ppr}}
    \right).
  \label{eq:signed-star-upper-work}
\end{equation}
\end{theorem}

\begin{proof}
The logarithm of $\lambda$ obeys
\begin{equation}
  -\log\lambda
  =\log\frac{1+t}{1-t}
  =2\operatorname{arctanh}(t)
  \ge2t,
  \label{eq:signed-star-lambda-exponential}
\end{equation}
so $\lambda^k\le e^{-2tk}$.  Put $y=tk$.  Since $t^2\le1$,
\begin{equation}
  \mathcal E(p^{(k)})
  \le\frac{(1+y)e^{-2y}}{m}.
  \label{eq:signed-star-error-envelope}
\end{equation}
For $k=k_\star$, $y\ge2$.  The function $(1+y)e^{-2y}$ is decreasing for
$y\ge0$, and
\begin{equation}
  (1+y)e^{-2y}\le3e^{-4}<\frac1{16}.
  \label{eq:signed-star-numerical-constant}
\end{equation}
By \eqref{eq:signed-star-size-bounds}, $m\ge1/(16\varepsilon_{\rm ppr})$;
hence $1/(16m)\le\varepsilon_{\rm ppr}$.  This proves accuracy.

Every block costs $m$.  The algorithm performs $k_\star$ blocks, reads one
seed entry, and writes $m+1$ output coordinates.  The star validation scans
can be combined with the first center and leaf blocks.  This proves the first
work inequality; $m\le1/(8\varepsilon_{\rm ppr})$ and
$\lceil2/t\rceil\le2/t+1$ give the asymptotic statement.
\end{proof}

\subsection{Why certificate stopping pays an avoidable logarithm}

At a block boundary, the largest degree-normalized residual magnitude is
\begin{equation}
  \max_u\frac{|r_u|}{d_u}
  =\frac{a_k}{m}
  =\frac{\gamma_\alpha(k+1)\lambda^k}{m}.
  \label{eq:signed-star-residual-at-block}
\end{equation}
The sufficient certificate \eqref{eq:signed-star-certificate} therefore
requires
\begin{equation}
  (k+1)\lambda^k\le\varepsilon_{\rm ppr}m.
  \label{eq:signed-star-certificate-block-condition}
\end{equation}
For our star, the right-hand side is a fixed constant at most $1/8$.
Because $-\log\lambda=\Theta(t)$, the descending solution of
\eqref{eq:signed-star-certificate-block-condition} is
\begin{equation}
  k=\Theta\!\left(\frac{\log(1/t)}{t}\right)
   =\Theta\!\left(
      \frac{\log(1/\alpha)}{\sqrt\alpha}
     \right).
  \label{eq:signed-star-certificate-log}
\end{equation}
By contrast, the semantic identity needs only $k=O(1/t)$.  Thus the
logarithm is a property of the stronger certificate on this signed trajectory,
not a lower bound on the algorithm's semantic complexity.  This distinction
is safe here because the predetermined block count proves accuracy without
querying the unknown solution.  On a general graph, replacing the certificate
requires a new locally checkable argument; the star formula cannot simply be
assumed.
